\partnote{Resolution, repair, and retained detail}{Resolution gives a concrete family of views. Its compatibility differential connects exact maps, protected information, and residual coordinates. Recovering every finite coordinate is then separated from convergence in a chosen norm.}
\section{Resolution towers and exact transport}\label{sec:towers}
\subsection{The sequential model}
A tower in a category \(\mathcal C\) is a sequence
\[
 \cdots\xrightarrow{\rho_2}X_2\xrightarrow{\rho_1}X_1
 \xrightarrow{\rho_0}X_0.
\]
For sets, a point is a sequence \(x_n\in X_n\) with \(\rho_n(x_{n+1})=x_n\). Composite restrictions turn the tower into a view system on \(\N\), with an arrow \(i\to j\) when \(i\le j\). Its section space is the same coherent point type, with the same coordinates.

A candidate transport from \(X\) to \(Y\) consists of maps \(f_n:X_n\to Y_n\). It is exact when
\begin{equation}\label{eq:exacttower}
 \rho^Y_n f_{n+1}=f_n\rho^X_n\quad\text{for all }n.
\end{equation}

\begin{proposition}[The exact tower category]\label{KS16}
Identity maps and composites of exact transports are exact. An exact transport induces a map of coherent point spaces by \(x\mapsto(f_nx_n)_n\), preserving identities and composition.
\end{proposition}
\begin{proof}
For two exact transports, substitute \(\rho^Z_ng_{n+1}=g_n\rho^Y_n\) and \(\rho^Y_nf_{n+1}=f_n\rho^X_n\) into the composite square. For a coherent point,
\[
 \rho^Y_n f_{n+1}x_{n+1}
 =f_n\rho^X_nx_{n+1}=f_nx_n.
\]
Identity, associativity, and composition on points are checked coordinatewise.
\end{proof}
\source{KS16}
This is a category on a fixed schedule. A full procategory allows additional reindexing and equivalence relations on maps. Those constructions must be supplied rather than identified with the fixed schedule by terminology.

\subsection{The actual residual of a transport}
Assume \(\mathcal C\) is preadditive, so differences of maps make sense. The tower differential is
\[
 (\partial f)_n=\rho^Y_nf_{n+1}-f_n\rho^X_n.
\]
Its target has components \(\Hom(X_{n+1},Y_n)\), not \(\Hom(X_n,Y_n)\).

\begin{proposition}[The transport Leibniz identity]\label{KS17}
For composable candidates,
\[
 \partial(gf)_n=(\partial g)_nf_{n+1}+g_n(\partial f)_n.
\]
In particular exactness is equivalent to vanishing of the differential.
\end{proposition}
\begin{proof}
Expand the right side. The terms \(-g_n\rho^Y_nf_{n+1}\) and \(+g_n\rho^Y_nf_{n+1}\) cancel, leaving \(\rho^Z_ng_{n+1}f_{n+1}-g_nf_n\rho^X_n\). This is the differential of the composite. The exactness statement is \eqref{eq:exacttower}.
\end{proof}
\source{KS17}
There is no additional mixed defect term in this identity. Such a term would require a different definition of the defect or a different expansion. Typed cancellation is a useful safeguard against importing an attractive but incorrect error formula.

\subsection{Cofinal views preserve points, not every measurement}
Let \(\phi:\N\to\N\) be monotone and unbounded. The sampled tower uses levels \(X_{\phi(n)}\) and their composite restrictions.

\begin{theorem}[Cofinal reconstruction of points]\label{KS18}
Sampling induces an equivalence between the coherent points of \(X\) and of its sampled tower.
\end{theorem}
\begin{proof}
Write $\rho_{i,j}:X_j\to X_i$ for the composite restriction when $i\le j$, with $\rho_{i,i}=I$. These maps satisfy $\rho_{i,j}\rho_{j,k}=\rho_{i,k}$. A coherent sampled point obeys
\[
 \rho_{\phi(n),\phi(m)}y_m=y_n\qquad(n\le m),
\]
by iterating its adjacent coherence equations.

For each $k$, unboundedness gives an index $n$ with $k\le\phi(n)$. Set $x_k=\rho_{k,\phi(n)}y_n$. If $m$ is another permissible index, choose $\ell\ge n,m$. Monotonicity permits both composite restrictions, and sampled coherence gives
\[
 \rho_{k,\phi(n)}y_n
 =\rho_{k,\phi(n)}\rho_{\phi(n),\phi(\ell)}y_\ell
 =\rho_{k,\phi(\ell)}y_\ell
 =\rho_{k,\phi(m)}y_m.
\]
Thus $x_k$ is independent of the choice, including when the schedule repeats a level.

For the adjacent equation at $k$, choose one $n$ with $k+1\le\phi(n)$. Then
\[
 \rho_{k,k+1}x_{k+1}
 =\rho_{k,k+1}\rho_{k+1,\phi(n)}y_n
 =\rho_{k,\phi(n)}y_n=x_k.
\]
The recovered family is therefore a point of the original tower.

Sampling a full point and recovering its $k$th coordinate gives $\rho_{k,\phi(n)}x_{\phi(n)}=x_k$ by its original coherence. Recovering a sampled point and evaluating at $k=\phi(n)$ allows the index $n$ itself, giving $\rho_{\phi(n),\phi(n)}y_n=y_n$. These are both inverse identities. No surjectivity of the bonding maps or of the point evaluations was used.
\end{proof}
\source{KS18}
This theorem is about exact coherent points. Raw residual size is a different quantity. On identity scalar towers, \(f_n=(-1)^n\) has nonzero fine differences, but its even subsequence is constant. The coarse residual vanishes because adjacent fine defects cancel. A cofinal equivalence of limits does not imply equality of residual norms or an absence of aliasing in uncorrected data.

\plate{03_tower}{Exact transport commutes with every adjacent restriction. A cofinal subsequence retains coherent points, but a raw defect can cancel during coarsening. The retained coarse value alone need not record the omitted fine discrepancy.}{fig:tower}

\chap{Finite observations and realized transformations}\label{sec:observation}
The point-space construction has a converse question. Given a function between coherent realization spaces, can it be represented by compatible maps between finite views? The answer is not automatically yes. This section constructs the relevant category of reindexed maps and identifies its realized functions under an explicit observability hypothesis.

\subsection{An output level may require a finer input level}
Let \(X\) and \(Y\) be sequential set-valued resolutions. Write \(\rho^X_{i,j}:X_j\to X_i\) for composite restriction. A reindexed transport consists of a monotone schedule \(\phi:\N\to\N\) and maps
\[
 f_n:X_{\phi(n)}\longrightarrow Y_n
\]
satisfying, for \(i\le j\),
\begin{equation}\label{eq:scheduled-naturality}
 \rho^Y_{i,j}f_j=f_i\rho^X_{\phi(i),\phi(j)}.
\end{equation}
The schedule records how much input is used for each output level. Monotonicity means that asking for more output never forces the representation to discard input already available. It does not require strict growth, a bounded look-ahead, or an efficiency estimate.

The realized map is
\[
 |f|:\operatorname{Point}(X)\to\operatorname{Point}(Y),\qquad
 |f|(x)_n=f_n(x_{\phi(n)}).
\]
Equation~\eqref{eq:scheduled-naturality} proves that this is a coherent output. A second transport \(g_n:Y_{\psi(n)}\to Z_n\) composes with \(f\) using schedule \(\phi\circ\psi\) and coordinate map \(g_n f_{\psi(n)}\). The order of the schedules is fixed by the types, not by the order in which their letters are written in prose.

\subsection{Equality is witnessed at finite stages}
Two different schedules can represent the same transformation. To compare \((\phi,f)\) and \((\psi,g)\), require that for every output level \(n\) there be a finite level \(k\ge\phi(n),\psi(n)\) with
\begin{equation}\label{eq:eventual-equality}
 f_n\rho^X_{\phi(n),k}=g_n\rho^X_{\psi(n),k}
 \quad\text{as functions on all of }X_k.
\end{equation}
Call this eventual equality. The witnessing \(k\) may depend on \(n\). Equality is required on the whole finite-level fiber, not only on states that happen to extend globally.

\begin{theorem}[The eventual transport category]\label{KS61}
Eventual equality is an equivalence relation and is preserved by composition. Its quotient therefore gives a category of resolutions and eventual transports. Realization descends to a functor into sets, preserving identities and composition.
\end{theorem}
\begin{proof}
A finite-stage equality remains true after moving to a later stage. Indeed, if it holds on $X_k$ and $\ell\ge k$, precompose both sides with $\rho^X_{k,\ell}$ and use the composite-restriction laws. Reflexivity uses $k=\phi(n)$, symmetry reverses the equality, and transitivity moves two witnesses to their maximum. The shared middle presentation then cancels as a literal function on that common fiber. This proves an equivalence relation without referring to global points.

We check compatibility with composition with all indices displayed. Write $f,f'$ with schedules $\phi,\phi'$ and $g,g'$ with schedules $\psi,\psi'$. Fix an output level $n$. Choose $j\ge\psi(n),\psi'(n)$ with
\[
 g_n\rho^Y_{\psi(n),j}=g'_n\rho^Y_{\psi'(n),j}.
\]
Next choose $k\ge\phi(j),\phi'(j)$ with
\[
 f_j\rho^X_{\phi(j),k}=f'_j\rho^X_{\phi'(j),k}.
\]
Monotonicity implies that $k$ is also at least both composite input levels. Applying the naturality equations and then the two witness equations gives
\[
\begin{aligned}
 g_nf_{\psi(n)}\rho^X_{\phi(\psi(n)),k}
 &=g_n\rho^Y_{\psi(n),j}f_j\rho^X_{\phi(j),k}\\
 &=g'_n\rho^Y_{\psi'(n),j}f_j\rho^X_{\phi(j),k}\\
 &=g'_n\rho^Y_{\psi'(n),j}f'_j\rho^X_{\phi'(j),k}\\
 &=g'_nf'_{\psi'(n)}\rho^X_{\phi'(\psi'(n)),k}.
\end{aligned}
\]
This is the required equality on every element of $X_k$. Thus composition descends to classes. The representative composition is associative because function composition and schedule composition are associative. The schedule $n\mapsto n$ with identity coordinate maps gives both identity laws. Hence the quotient has the claimed category structure.

Finally, evaluate a witness equation at the $k$th coordinate of a coherent point $x$. Its two restricted arguments become $x_{\phi(n)}$ and $x_{\psi(n)}$, proving equality of the realized outputs. Thus realization is well defined on the quotient. For a composite, direct evaluation gives
\[
 |g\circ f|(x)_n=g_n\bigl(f_{\psi(n)}(x_{\phi(\psi(n))})\bigr)
 =(|g|(|f|(x)))_n.
\]
The identity law is equally coordinatewise. Joint determination of output points proves the functor laws.
\end{proof}
\source{KS61}
This is an actual quotient category in the formal source. No identification with the full procategory of an arbitrary base category is needed or claimed. The objects here are sequential set-valued resolutions, and the equivalence relation is the finite-stage relation just defined.

\subsection{When global behavior determines the finite presentation}
Say that \(X\) has surjective point projections when
\[
 \forall k\quad \operatorname{ev}_k:\operatorname{Point}(X)\to X_k
 \quad\text{is onto}.
\]
This means that every finite input state really occurs in some coherent global input. It is stronger than nonemptiness of the global space.

\begin{theorem}[Faithfulness on an observable source]\label{KS62}
If \(X\) has surjective point projections, two reindexed maps out of \(X\) are eventually equal exactly when they induce the same realized function. Surjective adjacent bonding maps are a sufficient condition for surjective point projections.
\end{theorem}
\begin{proof}
Realized equality follows from eventual equality by the preceding theorem, so only its converse remains. Suppose \(|f|=|g|\). Fix the output level \(n\), and let \(k=\max\{\phi(n),\psi(n)\}\). To prove equality of two maps defined on \(X_k\), fix an arbitrary element \(a\in X_k\). The assumed surjectivity of \(\operatorname{ev}_k\) supplies a coherent point \(x\) with \(x_k=a\). Equality of the two realized maps gives equality of their output coordinates:
\[
 f_n(x_{\phi(n)})=g_n(x_{\psi(n)}).
\]
Coherence of \(x\) rewrites this as
\[
 f_n\bigl(\rho^X_{\phi(n),k}(a)\bigr)
 =g_n\bigl(\rho^X_{\psi(n),k}(a)\bigr).
\]
The index \(k\) was selected before \(a\), so this is one common finite-stage equality for \emph{all} elements of \(X_k\). It is exactly the witness required by eventual equality at output \(n\). Repeating for each \(n\) proves faithfulness.

For the sufficient condition on bonding maps, fix \(k\) and \(a\in X_k\). Define coordinates below \(k\) by restriction of \(a\). At each later stage use surjectivity of the adjacent bonding map to choose a preimage of the coordinate already chosen. Recursion with classical choice constructs the upper tail. Adjacent equations hold by construction, and induction on the interval length gives every composite equation. This produces a point evaluating to \(a\). If a fiber is empty, the statement that each of its elements lifts is vacuous, as required by surjectivity onto that fiber. Without the projection-surjectivity assumption, equality on realizable inputs alone need not identify finite maps on values that never occur in a global point.
\end{proof}
\source{KS62}
The hypothesis is sufficient, not presented as a characterization of every source on which faithfulness happens to hold. Its role in the proof is precise: it converts a statement about globally realized inputs into a statement about all finite inputs.

Without it, global behavior can lose information. For example, use \(X_n=\{m\in\N:m\ge n\}\) with inclusion bonds and a constant Boolean output tower. The source point space is empty. The two transports that always output false and always output true induce the same function from the empty set, but disagree at every finite stage. Every \(X_k\) is nonempty, so no eventual equality witness exists. This calculation uses the finite-prefix boundary of KS54 and does not identify empty global behavior with equality of finite presentations.

\subsection{The exact image of realization}
For a function \(F:\operatorname{Point}(X)\to\operatorname{Point}(Y)\), define finite observation by the quantifiers
\begin{equation}\label{eq:finite-observation}
 \forall n\ \exists k\ \exists a:X_k\to Y_n\ \forall x:\quad
 F(x)_n=a(x_k).
\end{equation}
The input level is uniform over all global inputs for each fixed output level. It may vary with the output level. This is not the weaker condition that each individual input can be inspected to some input-dependent finite depth.

\begin{theorem}[Characterization of realized transformations]\label{KS63}
If \(X\) has surjective point projections, realization induces a bijection
\[
 \operatorname{EventualHom}(X,Y)
 \simeq
 \{F:\operatorname{Point}(X)\to\operatorname{Point}(Y):F
 \text{ is finitely observed}\}.
\]
\end{theorem}
\begin{proof}
For a reindexed map $(\phi,f)$, the formula $|f|(x)_n=f_n(x_{\phi(n)})$ immediately supplies a finite-observation witness for each $n$. Eventual equality does not affect the realized function by KS61, so this defines a map from the quotient to the stated set of functions.

For surjectivity, let $F$ be finitely observed. Choose one pair $(k_n,a_n)$ at each output level, satisfying $F(x)_n=a_n(x_{k_n})$ for every coherent input. Define
\[
 \phi(n)=\max_{0\le i\le n}k_i,
 \qquad f_n=a_n\rho^X_{k_n,\phi(n)}.
\]
The maximum is of a nonempty finite set. It gives a monotone schedule with $k_n\le\phi(n)$. For every point $x$, coherence in $X$ gives
\[
 f_n(x_{\phi(n)})
 =a_n\bigl(\rho^X_{k_n,\phi(n)}x_{\phi(n)}\bigr)
 =a_n(x_{k_n})=F(x)_n.                                      \tag{*}
\]

The remaining obligation is naturality on \emph{all} finite states. Fix $i\le j$ and $u\in X_{\phi(j)}$. By surjectivity of that point projection, choose $x$ with $x_{\phi(j)}=u$. Equation (*) and output coherence yield
\[
\begin{aligned}
 \rho^Y_{i,j}f_j(u)&=\rho^Y_{i,j}F(x)_j=F(x)_i,\\
 f_i\rho^X_{\phi(i),\phi(j)}(u)&=f_i(x_{\phi(i)})=F(x)_i.
\end{aligned}
\]
The state $u$ was arbitrary. Hence the two finite-level functions are equal, which proves scheduled naturality. The resulting transport realizes $F$ by (*).

Different witness choices can give different schedules and coordinate maps, but the realized function is still $F$. By KS62 these transports are eventually equal, so their quotient class is independent of all choices. The same faithfulness result says that two classes with equal realized functions are identical. We have proved both surjectivity and injectivity, including the well-definedness of the inverse choice construction. No bound on the growth of $\phi$ follows from these quantifiers.
\end{proof}
\source{KS63}
The finite-level compatibility is derived in this proof. It is not inserted as an additional field in the hypothesis that \(F\) is finitely observed. This is one of the useful bridges between a global description of a transformation and an actual presentation in the resolution calculus.

\fixedplate{20_schedule}{Independent finite-observation witnesses need not use a monotone input schedule. Taking the running maximum produces one without changing the realized function. The map at the enlarged level first restricts to the originally sufficient level. This construction gives no efficiency bound on how large those levels are.}{fig:schedule}

\subsection{Binary prefixes: a complete interpretation}
Let \(X_n\) consist of binary words of length \(n\), with prefix restriction. A coherent point is exactly an infinite binary sequence: its length-\(n\) view is its first \(n\) bits. Every prefix extends, for example by appending zeros. Hence the source satisfies the hypothesis of the characterization theorem.

Bitwise complementation is represented at level \(n\) by complementing the first \(n\) bits. The left shift \(F(x)_i=x_{i+1}\) uses the schedule \(\phi(n)=n+1\). A nearest-neighbor difference over \(\mathbb F_2\),
\[
 G(x)_i=x_i+x_{i+1},
\]
also uses \(n+1\) input bits for \(n\) output bits. Applying \(G\) twice uses \(n+2\) bits, and the middle terms cancel:
\[
 G(G(x))_i=x_i+x_{i+2}.
\]
The compatible finite maps make each statement meaningful on the entire coherent sequence, not merely on one chosen truncation. These are worked instances of the scheduled formulas. The finite regression suite checks their prefix compatibility independently.

Not every set function on these points is finitely observed. Consider the Boolean output that is one exactly when every input bit is zero, regarded as a point of the constant Boolean tower. Given any proposed finite level \(k\), compare the all-zero sequence with a sequence whose first one occurs at position \(k\). Their first \(k\) bits agree, but the outputs differ. Therefore no finite witness \(a:X_k\to\{0,1\}\) can compute even that one output coordinate. This rules out a reindexed presentation by the forward direction of the theorem.

The obstruction in this example is one of information, not a runtime benchmark. No amount of computation on the specified finite prefix distinguishes those two inputs. The theorem does not prohibit a different mathematical representation that supplies more information as part of each view. Changing the available view is changing the input model.

\subsection{Cofinal changes also preserve the morphism calculus}
A monotone unbounded schedule \(\theta\) gives the sampled resolution \(X_\theta\). There is a transport from \(X\) to \(X_\theta\) that reads the selected level. For the reverse direction choose, for each \(n\), a level \(j(n)\) with \(n\le\theta(j(n))\), and restrict the available value down to \(X_n\). Taking a least such index gives a monotone schedule.

\begin{theorem}[Cofinal invariance of eventual transports]\label{KS64}
Sampling and recovery are inverse in the eventual transport category. Consequently, for monotone unbounded schedules \(\theta\) and \(\eta\), there is a bijection
\[
 \operatorname{EventualHom}(X,Y)
 \simeq
 \operatorname{EventualHom}(X_\theta,Y_\eta).
\]
On point spaces it agrees with the previously constructed cofinal realization equivalences.
\end{theorem}
\begin{proof}
Let \(s:X\to X_\theta\) be sampling, with schedule \(\theta\) and identity level maps. For recovery choose the least \(j(n)\) with \(n\le\theta(j(n))\). Unboundedness makes this set nonempty. Increasing \(n\) only strengthens the inequality, so \(j\) is monotone. The level map of \(r:X_\theta\to X\) is \(\rho^X_{n,\theta(j(n))}\). Composite restriction proves its naturality.

The composite \(rs:X\to X\) reads level \(\theta(j(n))\) and restricts to \(n\). At that very stage its map agrees with the identity transport after restriction from \(\theta(j(n))\) to \(n\). Thus \(rs\) is eventually equal to the identity.

For \(sr:X_\theta\to X_\theta\), the schedule at output \(n\) is \(j(\theta(n))\), and the level map is \(\rho^X_{\theta(n),\theta(j(\theta(n)))}\). Compare it with the identity at a sampled index
\[
 k\ge\max\{n,j(\theta(n))\}.
\]
Restricting either route from level \(k\) of the sampled system gives the same original restriction \(\rho^X_{\theta(n),\theta(k)}\). This proves eventual equality even when \(\theta\) has repeated values. Hence the two classes are inverse isomorphisms, without requiring them to be literal inverse schedules.

Write \(s_X,r_X\) and \(s_Y,r_Y\) for the sampling and recovery isomorphisms of source and target. The asserted Hom map is
\[
 [f]\longmapsto s_Y[f]r_X,
\]
with inverse \([g]\mapsto r_Y[g]s_X\). Associativity and the two inverse equations cancel the inserted pairs, proving both round trips. Applying the realization functor transports these equations to point spaces, where sampling and recovery are the cofinal point equivalences already constructed. No norm invariance or surjectivity of source point projections is needed for this categorical conjugation argument.
\end{proof}
\source{KS64}
This strengthens cofinal invariance from points to a specified category of transformations. It still does not make raw residual norms invariant, identify every global set function with a finite presentation, or preserve the numerical conditioning of a repair. Those are separate questions, addressed by the defect maps in the next sections.

\Needspace{7\baselineskip}
\begin{insight}
For a projection-surjective source, eventual finite presentations correspond exactly to uniformly finite observations at each output level. This is a precise machine-oriented representation theorem. It is not a claim that every represented map is computationally efficient, or that every arbitrary function of an infinite input has such a presentation.
\end{insight}


\chap{Banach repair and constructed split coordinates}\label{sec:split}
\subsection{An actual complete candidate space}
Let \(X_n,Y_n\) be Banach spaces over a complete nontrivially normed field, with continuous linear bonding maps and uniform bounds \(\norm{\rho^X_n}\le B_X\), \(\norm{\rho^Y_n}\le B_Y\). Use the spaces of bounded operator families
\[
 E=\ell^\infty_n\mathcal L(X_n,Y_n),\qquad
 R=\ell^\infty_n\mathcal L(X_{n+1},Y_n).
\]
The source formalization derives the required completeness from the operator spaces and their bounded product. The same differential satisfies
\[
 \norm{\partial f}_\infty\le(B_X+B_Y)\norm f_\infty.
\]
Thus it is a continuous linear map, and evaluation at a fixed level is continuous. These facts connect algebraic compatibility to an analytic limit argument.

\subsection{From a correction operator to a projection}
The next statement is general. Let \(E\) be Banach, let \(d:E\to R\), \(p:E\to S\), and \(A:R\to E\) be continuous linear maps. Suppose \(pA=0\) and, for \(0\le q<1\),
\begin{equation}\label{eq:contraction}
 \norm{df-dA(df)}\le q\norm{df}\qquad(f\in E).
\end{equation}
The hypothesis concerns residuals of actual candidates. It need not hold for every independent vector in \(R\).

\begin{theorem}[Constructed geometric repair]\label{KS19}
Starting from \(f^{(0)}=f\), define
\[
 f^{(m+1)}=f^{(m)}-A(df^{(m)}).
\]
There is a bounded linear projection \(P_A:E\to\ker d\) with
\[
 f^{(m)}\longrightarrow P_Af,\qquad p(P_Af)=p(f),
\]
and
\[
 \norm{f^{(m)}-P_Af}
 \le\frac{\norm A}{1-q}\norm{df^{(m)}}.
\]
It fixes every exact candidate and satisfies
\[
 \norm{P_A}\le1+\frac{\norm A\norm d}{1-q}.
\]
\end{theorem}
\begin{proof}
Put $r_m=df^{(m)}$. Additivity of $d$ gives
\[
 r_{m+1}=r_m-dAr_m.
\]
Each $r_m$ is the residual of an actual candidate, so the stated hypothesis applies at every iteration. Induction yields, for $j\ge0$,
\[
 \norm{r_{m+j}}\le q^j\norm{r_m},\qquad
 \norm{f^{(m+j+1)}-f^{(m+j)}}
 \le\norm A\,q^j\norm{r_m}.
\]
For $n>m$, the finite triangle inequality therefore gives
\[
 \norm{f^{(n)}-f^{(m)}}
 \le\norm A\,\norm{r_m}\sum_{j=0}^{n-m-1}q^j
 \le\frac{\norm A}{1-q}\norm{r_m}.
\]
Since $\norm{r_m}\le q^m\norm{df}$ tends to zero, this is a Cauchy estimate. Completeness of $E$ supplies a limit, denoted $P_Af$. Passing $n\to\infty$ in the same norm inequality proves the required tail estimate.

Continuity of $d$ implies $d(P_Af)=\lim_m r_m=0$. Protection holds before taking a limit: $pA=0$ gives $pf^{(m+1)}=pf^{(m)}$, hence $pf^{(m)}=pf$ for every $m$. Continuity of $p$ gives $p(P_Af)=pf$. This explains why continuity of the observable belongs to this analytic statement.

For linearity, write $L=I-Ad$, so $f^{(m)}=L^mf$. For fixed $f,g$ and scalars $\lambda,\mu$,
\[
 L^m(\lambda f+\mu g)=\lambda L^mf+\mu L^mg.
\]
Continuity of addition and scalar multiplication allows passage to the norm limit and proves linearity of $P_A$. The estimate with $m=0$ gives
\[
 \norm{P_Af}\le\norm f+\frac{\norm A}{1-q}\norm{df}
 \le\left(1+\frac{\norm A\norm d}{1-q}\right)\norm f,
\]
so the linear map is bounded. If $df=0$, every iterate equals $f$, hence $P_Af=f$. Its image is contained in $\ker d$ and it fixes that kernel, proving surjectivity onto its stated image and $P_A^2=P_A$. For $q=0$ the residual vanishes after the first update and the same inequalities remain valid.
\end{proof}
\source{KS19}
This theorem constructs the dynamics from \(A\). It does not construct a suitable \(A\) for every residual problem. That missing input can, however, be derived for a useful class of actual towers.

\subsection{A right inverse from target splittings}
Suppose each target bonding map has a continuous linear section
\[
 s_n:Y_n\to Y_{n+1},\qquad \rho^Y_ns_n=I,\qquad \norm{s_n}\le S.
\]
Assume the propagation gap
\begin{equation}\label{eq:splitgap}
 SB_X<1,\qquad K=\frac{S}{1-SB_X}.
\end{equation}
For \(r\in R\), recursively set
\begin{equation}\label{eq:greenrecurrence}
 (Ar)_0=0,\qquad
 (Ar)_{n+1}=s_n\bigl(r_n+(Ar)_n\rho^X_n\bigr).
\end{equation}

\begin{theorem}[Split tower correction]\label{KS20}
The recurrence defines a continuous linear map \(A:R\to E\) with
\[
 \partial A=I_R,\qquad \norm A\le K.
\]
The operator \(P=I-A\partial\) is an exact projection and
\[
 (Pf)_0=f_0,\qquad
 \norm{f-Pf}_\infty\le K\norm{\partial f}_\infty.
\]
\end{theorem}
\begin{proof}
Let $\theta=SB_X<1$. The number $K=S/(1-\theta)$ is nonnegative and satisfies
\[
 S+SB_XK=K.
\]
Fix $r\in R$ and define its coordinates recursively as in \eqref{eq:greenrecurrence}. The initial coordinate has norm zero. If $\norm{(Ar)_n}\le K\norm r_\infty$, submultiplicativity of operator norms gives
\[
 \begin{aligned}
 \norm{(Ar)_{n+1}}
 &\le\norm{s_n}\bigl(\norm{r_n}+\norm{(Ar)_n}\norm{\rho^X_n}\bigr)\\
 &\le S(1+KB_X)\norm r_\infty=K\norm r_\infty.
 \end{aligned}
\]
Induction proves the estimate at every coordinate, so the recursively defined family actually belongs to the bounded product $E$.

The coordinate recurrence preserves addition and scalar multiplication, starting from the zero initial coordinate. Induction therefore proves these laws for the family map $A$. Its uniform estimate proves $\norm A\le K$, giving a continuous linear map with the claimed domain and codomain.

For every $n$, apply the target restriction to the recurrence. The equation $\rho^Y_ns_n=I$ gives
\[
 \rho^Y_n(Ar)_{n+1}=r_n+(Ar)_n\rho^X_n,
 \qquad (\partial Ar)_n=r_n.
\]
Coordinate equality yields the operator identity $\partial A=I_R$. In particular this identity holds for every bounded residual, not just residuals already known to lie in the image.

Set $P=I-A\partial$. Then
\[
 \partial P=\partial-\partial A\partial=0,\qquad
 P^2=I-2A\partial+A(\partial A)\partial=I-A\partial=P.
\]
The expression $2A\partial$ abbreviates addition of the same operator twice and uses no division by two. Every exact $f$ is fixed because $\partial f=0$. Since $(Ar)_0=0$, the seed of $Pf$ is $f_0$. Finally $f-Pf=A\partial f$, and the bound on $A$ gives the asserted residual estimate. The denominators above are justified exactly by the strict propagation gap.
\end{proof}
\source{KS20}
The surjectivity is onto all bounded residual families. It is stronger than the assertion that some exact output can be assigned to each candidate. The gap and the chosen splittings are genuine inputs.

\subsection{Protected information beyond the seed}
When the target grows, its new kernel directions are independent information. Define
\[
 N_n(f)=(I-s_n\rho^Y_n)f_{n+1},\qquad
 (p_Lf)_0=f_0,\quad(p_Lf)_{n+1}=N_n(f).
\]
Each \(N_n(f)\) lands in \(\ker\rho^Y_n\). The observable \(p_L\) retains both the seed and every new output component, relative to the chosen sections.

\begin{theorem}[Split coordinate theorem]\label{KS21}
The observable is a bounded projection with \(p_LA=0\) and \(p_LP=p_L\). It separates exact candidates. If
\[
 \mathcal S_L=\{s\in E:p_Ls=s\},
\]
then
\[
 \Phi:E\simeq_L\mathcal S_L\times R,
 \qquad \Phi(f)=(p_Lf,\partial f),
 \qquad \Phi^{-1}(s,r)=Ps+Ar.
\]
Repair becomes \((s,r)\mapsto(s,0)\). It is the unique exact repair preserving \(p_Lf\).
\end{theorem}
\begin{proof}
Put \(B_n=I-s_n\rho^Y_n\). Since \(\rho^Y_ns_n=I\), direct multiplication gives
\[
 B_n^2=B_n,\qquad \rho^Y_nB_n=0,\qquad B_ns_n=0.
\]
Moreover, \(B_ny=y\) whenever \(\rho^Y_ny=0\). Thus \(B_n\) really projects onto the target kernel. Its norm is at most \(1+SB_Y\). Coordinate zero of \(p_L\) is the identity, and every later coordinate applies \(B_n\), so \(p_L^2=p_L\) and \(\norm{p_L}\le1+SB_Y\). In particular the formula defines a map on the bounded product, not just on arbitrary coordinate families.

The recurrence for \(Ar\) gives zero seed and places \((Ar)_{n+1}\) in the image of \(s_n\). The identities just proved make all its innovations zero. Therefore \(p_LA=0\), and
\[
 p_LP=p_L(I-A\partial)=p_L.
\]
We next prove separation on the exact space. If \(f\) is exact, its restriction equation rewrites as
\[
 f_{n+1}=s_n\rho^Y_nf_{n+1}+B_nf_{n+1}
          =s_nf_n\rho^X_n+N_n(f).
\]
Two exact families with the same seed and innovations agree first at level zero and then at every later level by this recurrence. More generally, if two arbitrary families have the same observable and differential, their difference is exact with zero observable, so separation on the exact space makes their difference zero.

The proposed chart has the claimed target because \(p_L(p_Lf)=p_Lf\). Given \(s\in\mathcal S_L\) and \(r\in R\), set \(f=Ps+Ar\). The preceding identities and \(\partial A=I\) give
\[
 p_Lf=p_Ls=s,\qquad \partial f=r.
\]
This proves one round trip. For the other, \(f\) and \(P(p_Lf)+A\partial f\) have the same observable and residual, so the joint separation just proved identifies them. All maps are linear and bounded on their indicated subspaces and finite products, establishing a continuous linear equivalence.

Finally, \(Pf\) has coordinates \((p_Lf,0)\). Any exact answer with the same observable has those same two coordinates, hence equals \(Pf\). This uniqueness is relative to the complete observable of seed and innovations. It is not obtained by assuming that the seed alone determines an exact growing-target family.
\end{proof}
\source{KS21}
This is an actual analytic coordinate equivalence, not generally an isometry. The protected coordinate depends on the declared splitting. It should not be called an intrinsic semantic coordinate without a further comparison theorem.

\fixedplate{04_split}{In the split coordinate chart, a candidate separates into retained information and compatibility residual. The exact subspace is the zero residual slice. The chart and its protected coordinate depend on the chosen target sections.}{fig:split}

\subsection{A growing example and the gap boundary}
Use \(X_n=Y_n=\R^{n+1}\) with maximum norms. Let \(\pi_n\) delete the last coordinate, take \(\rho^X_n=\tfrac12\pi_n\), \(\rho^Y_n=\pi_n\), and let \(s_n\) append zero. Then \(S=B_Y=1\), \(B_X=1/2\), and \(K=2\).

\begin{example}[Repair of the identity family]\label{KS22}
The bounded candidate \(f_n=I\) is not exact. Its repair is
\[
 (Pf)_n=\operatorname{diag}(2^{-n},2^{-(n-1)},\ldots,2^{-1},1).
\]
The new last output row is retained at every stage while inherited rows are propagated by the factor \(1/2\).
\end{example}
\begin{proof}
At level zero the value is one. Exactness forces every inherited row at the next level to be one half of its previous value. Preservation of the innovation forces the new row to remain the corresponding identity row. Induction gives the displayed diagonal.
\end{proof}
\source{KS22}
For constant real scalar identity towers, by contrast, \((\partial f)_n=f_{n+1}-f_n\). The constant residual \(r_n=1\) has no bounded primitive, since any primitive satisfies \(f_n=f_0+n\). Here \(SB_X=1\). This counterexample excludes discarding the gap indiscriminately. It does not prove that the gap is necessary for every individual tower. Exact repair preserving the initial coordinate still exists in this constant example, by replacing a candidate with its initial constant family. Surjectivity onto all bounded residuals does not follow from that repair.
\source{KS23}

\subsection{What is canonical}
For fixed splitting and protected observable, the repair above is unique. Changing the splitting can change both the observable and the repaired answer. In the growing example, replace \(s_n(y)=(y,0)\) by \(s_n'(y)=(y,y_0)\). At the first positive level, repair of the identity becomes
\[
 \begin{pmatrix}1/2&0\\-1/2&1\end{pmatrix}
 \quad\text{instead of}\quad
 \begin{pmatrix}1/2&0\\0&1\end{pmatrix}.
\]
Both are exact continuations for their respective innovation profiles and have the same seed. Thus a unique answer for a declared preservation problem is not an answer independent of all preservation choices.
\source{KS24}

\chap{Repair under a change of representation}\label{sec:repair-transport}
An exact transformation should carry coherent inputs to coherent outputs. It need not carry the selected repair of an incoherent input to the selected repair of its transformed image. The distinction has an exact formula. In split coordinates, a compatible change of representation is generally block triangular rather than block diagonal. Its off-diagonal term records precisely the discrepancy between the two repair routes.

\subsection{A section is a choice of residual coordinates}
Let \(d:E\to R\) be a continuous linear map between normed spaces over a nontrivially normed field. Suppose it has a continuous linear right inverse \(A:R\to E\), so \(dA=I_R\). Set
\[
 K=\ker d,\qquad P_A=I_E-Ad.
\]
Then \(dP_A=0\), \(P_A\) fixes \(K\), and \(P_A^2=P_A\). The right inverse chooses where each residual is represented inside the candidate space. Different right inverses can choose different representatives of the same residual.

\begin{theorem}[The split defect chart]\label{KS65}
The maps
\[
 \chi_A:E\simeq_L K\times R,\qquad
 \chi_A(x)=(P_Ax,dx),\qquad
 \chi_A^{-1}(u,r)=u+Ar
\]
are inverse continuous linear equivalences. In these coordinates, \(P_A\) sends \((u,r)\) to \((u,0)\). If \(B\) is another right inverse, then \(A-B\) takes values in \(K\), and
\[
 P_Bx-P_Ax=(A-B)dx.
\]
The projections commute if and only if they are equal, and they are equal if and only if \(A=B\).
\end{theorem}
\begin{proof}
The identity \(dA=I\) gives \(dP_A=d-dAd=0\), so the proposed first chart coordinate actually belongs to \(K\). If \(u\in K\), then \(du=0\) and \(P_Au=u\). Therefore, for \((u,r)\in K\times R\),
\[
 d(u+Ar)=r,\qquad P_A(u+Ar)=u+Ar-A(r)=u.
\]
In the other direction, \(P_Ax+A(dx)=x-Adx+Adx=x\). These are both inverse equations, including the membership condition for the kernel coordinate. All component maps are bounded linear, so they define a continuous linear equivalence on the standard finite product. Applied to \(P_Ax\), the chart gives \((P_Ax,0)\), proving the repair formula.

For a second right inverse \(B\), the equation \(d(A-B)=I-I=0\) places its difference from \(A\) in \(K\). Direct subtraction gives
\[
 P_B-P_A=(I-Bd)-(I-Ad)=(A-B)d.
\]
Both projections have range \(K\) and fix every point of \(K\). Thus applying \(P_A\) after \(P_B\) gives \(P_B\), while applying \(P_B\) after \(P_A\) gives \(P_A\). These two compositions are equal if and only if the projections are equal.

If \(P_A=P_B\), evaluate their difference at \(Ar\). Since \(dAr=r\), the difference formula yields \((A-B)r=0\) for every \(r\), hence \(A=B\). Equality of sections conversely makes their projection formulas identical. The argument uses the right inverse to test every residual. It does not infer equality of sections merely from agreement on some restricted set of observed residuals.
\end{proof}
\source{KS65}
This is a linear splitting, not necessarily an orthogonal one. The exact coordinate \(u\) depends on the chosen section. The residual coordinate \(r=dx\) does not. A specified Hilbert geometry can make one section distinguished through orthogonality, but the algebraic split alone does not do so.

The tower construction is stronger in a different direction. Its recurrence builds a right inverse from the target splittings, and its observable identifies what that repair preserves. The abstract chart here isolates the mechanism shared by those concrete constructions.

\subsection{A change of representation has two components}
Let \(D:F\to T\) be a second continuous linear defect map. A compatible change consists of bounded linear maps
\[
 S:E\to F,\qquad C:R\to T,\qquad DS=Cd.
\]
The equation is the commuting defect square. It says that transforming a candidate and then measuring its defect agrees with transforming its actual defect. No sections have yet been used.

In particular, \(S\) carries \(\ker d\) into \(\ker D\). It may still have a kernel, lose information, or fail to be onto. Calling it a change of representation does not assert an equivalence unless inverse maps have also been supplied.

Choose sections \(A\) of \(d\) and \(B\) of \(D\). Define
\begin{equation}\label{eq:repair-mismatch}
 M=SA-BC:R\longrightarrow F.
\end{equation}
The commuting square and section identities give \(DM=0\). Thus this mismatch takes residual information into an exact direction of the target.

\begin{theorem}[Triangular transport and repair mismatch]\label{KS66}
In split defect coordinates, the change of representation is
\begin{equation}\label{eq:triangular-chart}
 \chi_BS\chi_A^{-1}(u,r)=(Su+Mr,Cr).
\end{equation}
The two repair routes differ by
\begin{equation}\label{eq:repair-change}
 P_BSx-SP_Ax=M(dx).
\end{equation}
Consequently they agree on exact inputs, and
\[
 \norm{P_BSx-SP_Ax}\le\norm M\norm{dx},
 \qquad \norm M\le\norm S\norm A+\norm B\norm C.
\]
They agree on all inputs if and only if \(SA=BC\).

For two consecutive compatible changes \((S_1,C_1)\), \((S_2,C_2)\), with compatible choices of sections at the three ends, their mismatches satisfy
\[
 M_{21}=S_2M_1+M_2C_1.
\]
If \(dx_j\to0\) in norm for a fixed change and fixed sections, the repair discrepancy tends to zero.
\end{theorem}
\begin{proof}
First \(DS=Cd\) sends \(\ker d\) into \(\ker D\), and
\[
 DM=DSA-DBC=CdA-C=C-C=0.
\]
Thus both terms in the proposed exact coordinate have the correct codomain. Given \(u\in\ker d\), compute the residual and exact coordinates of \(S(u+Ar)\) separately:
\[
 DS(u+Ar)=Cr,\qquad
 S(u+Ar)-BCr=Su+(SA-BC)r.
\]
This proves the triangular formula as an equality in \(\ker D\times T\).

For a general input, use the same commuting square to obtain
\[
 \begin{aligned}
 P_BSx-SP_Ax
 &=Sx-BDSx-Sx+SAdx\\
 &=(SA-BC)dx=M(dx).
 \end{aligned}
\]
Its norm is at most \(\norm M\norm{dx}\). The triangle inequality and submultiplicativity give \(\norm M\le\norm S\norm A+\norm B\norm C\). Exact inputs have \(dx=0\), so the two routes agree there regardless of the section choices. If they agree on \emph{every} input, substitute \(x=Ar\). Then \(dAr=r\) makes \(Mr=0\) for each residual, hence \(SA=BC\). That equality conversely makes the discrepancy vanish everywhere.

For composition, denote the three sections by \(A,B,A_2\). The composite residual map is \(C_2C_1\), so its mismatch is \(S_2S_1A-A_2C_2C_1\). Insert and subtract \(S_2BC_1\):
\[
 \begin{aligned}
 S_2S_1A-A_2C_2C_1
 &=S_2(S_1A-BC_1)+(S_2B-A_2C_2)C_1\\
 &=S_2M_1+M_2C_1.
 \end{aligned}
\]
All terms map the first residual space into the last candidate space, in the written order. Finally, for fixed bounded \(M\), the estimate \(\norm{M(dx_j)}\le\norm M\norm{dx_j}\) proves convergence of the discrepancy when \(dx_j\to0\). Allowing the section or representation norms to grow with \(j\) would require a separate uniform bound.
\end{proof}
\source{KS66}
The off-diagonal term in \eqref{eq:triangular-chart} has a specific meaning. A raw residual can become part of the target's chosen exact coordinate after a change of presentation. Deleting the residual before changing presentation removes that contribution. Deleting it afterwards does not. This is not a failure of exactness preservation. Both routes end in the target exact space, but they can end at different points in it.

\subsection{A two-coordinate model in which the routes differ}
Use \(E=F=\R^2\), write candidates as \((u,r)\), and take
\[
 d(u,r)=D(u,r)=r,\qquad A(r)=B(r)=(0,r).
\]
The exact space is the horizontal axis and both repairs are \(P(u,r)=(u,0)\). For a real parameter \(c\), set
\[
 S_c(u,r)=(u+cr,r),\qquad C=I.
\]
The defect square commutes because the second coordinate is unchanged. The mismatch is \(M_c(r)=(cr,0)\). Hence
\[
 PS_c(u,r)=(u+cr,0),\qquad S_cP(u,r)=(u,0).
\]
With \(c=2\) and input \((0,1)\), one route gives \((2,0)\) and the other gives \((0,0)\). Both have zero residual. A report containing only their final residuals would miss the entire difference.

\plate{21_mismatch}{The shear preserves the residual coordinate but mixes it into the chosen exact coordinate. Repair after the shear reaches \((2,0)\). Repair before the shear reaches \((0,0)\). Their difference is the exact-valued mismatch applied to the original residual. The example uses the ordinary Euclidean drawing, while the algebraic identity does not depend on that drawing.}{fig:mismatch}

Two shears compose by adding their parameters. In the mismatch law this appears as \(M_{c_2+c_1}=S_{c_2}M_{c_1}+M_{c_2}\), and \(S_{c_2}\) fixes the exact horizontal axis. The identity is therefore a simple instance of the same composition rule, not an unrelated coordinate trick.

\subsection{Actual tower resampling gives a defect map}
Return to Banach towers with contractive source and target bonds. Let \(\phi\) be a monotone schedule. Candidate sampling is
\[
 (S_\phi f)_n=f_{\phi(n)}.
\]
The sampled bonding map is the actual composite restriction between the selected levels. The residual map \(C_\phi\) is not pointwise sampling of a scalar error score. It accumulates transported fine defects across each omitted block.

For a block from \(a\) to \(b\), its value on a residual family \(r\) is
\begin{equation}\label{eq:block-defect-sum}
 (C_{a,b}r)=\sum_{k=a}^{b-1}
 \rho^Y_{a,k}\,r_k\,\rho^X_{k+1,b}.
\end{equation}
Every summand maps \(X_b\) into \(Y_a\). These surrounding restrictions are necessary when the spaces change with the level. For an actual differential, the sum telescopes to
\[
 C_{a,b}(\partial f)=\rho^Y_{a,b}f_b-f_a\rho^X_{a,b}.
\]
The dyadic coarsener in Section~\ref{sec:cascade} is the case \(b=a+2\).

\begin{theorem}[Bounded-gap resampling and selected repair]\label{KS67}
Suppose \(\phi(n+1)-\phi(n)\le m\) for a fixed natural number \(m\). Then the actual sampling maps are bounded and satisfy
\[
 \norm{S_\phi}\le1,\qquad \norm{C_\phi}\le m,
 \qquad \partial_\phi S_\phi=C_\phi\partial.
\]
Exact candidates remain exact under every monotone common resampling, even without a gap bound. If \(A\) and \(B\) are bounded right inverses for the fine and sampled differentials, then
\[
 \norm{P_BS_\phi f-S_\phi P_Af}
 \le(\norm A+m\norm B)\norm{\partial f}.
\]
\end{theorem}
\begin{proof}
For a bounded candidate family, \(\sup_n\norm{f_{\phi(n)}}\le\sup_k\norm{f_k}\), proving the sampling bound. For a bounded residual family, each summand in \eqref{eq:block-defect-sum} is a map from \(X_b\) to \(Y_a\) and has norm at most \(\norm r_\infty\), because every composite bond is contractive. A selected block has length at most \(m\), so the triangle inequality bounds its sum by \(m\norm r_\infty\). Taking the supremum over blocks proves that \(C_\phi\) maps into the bounded residual product and has the asserted operator norm. If a block has length zero its sum is zero, in agreement with the identity restriction.

To verify the commuting square, insert
\(r_k=\rho^Y_kf_{k+1}-f_k\rho^X_k\). The \(k\)-th term becomes
\[
 \rho^Y_{a,k+1}f_{k+1}\rho^X_{k+1,b}
 -\rho^Y_{a,k}f_k\rho^X_{k,b}.
\]
Successive internal terms cancel in the finite sum, leaving \(\rho^Y_{a,b}f_b-f_a\rho^X_{a,b}\). This is exactly the sampled differential. No infinite series interchange is involved.

If the fine differential vanishes, the finite telescope is zero on every selected interval. That conclusion remains true for unbounded gap lengths, because each individual interval is finite. What may fail without a gap bound is boundedness of the residual transformation on \emph{arbitrary} bounded residual families, not preservation of exact inputs.

When the two bounded sections exist, Theorem~\ref{KS66} applies to the commuting square just established. Substitute \(\norm{S_\phi}\le1\) and \(\norm{C_\phi}\le m\) in its estimate to get \((\norm A+m\norm B)\norm{\partial f}\). Existence of those sections is an additional hypothesis, not a consequence of the sampling estimate.
\end{proof}
\source{KS67}
The existence of the two right inverses is not inferred merely from contractivity of the bonds. The split-tower theorem provides them under its additional section and propagation-gap conditions. Keeping these hypotheses separate prevents a resampling theorem from being mistaken for an unconditional reconstruction algorithm.

\subsection{Two distinct losses under coarsening}
On scalar identity towers, a fine residual block \((1,-1)\) has zero two-step sum. Thus \(C_\phi r=0\) need not imply \(r=0\), even for a bounded-gap schedule. This is cancellation. Retaining an independent detail coordinate remedies that information loss in a split coarsener, as the next section shows.

Unbounded gaps cause a different issue. With \(\phi(n)=n^2\), the bounded residual family \(r_k=1\) has block sums \(2n+1\). These are unbounded. Therefore block accumulation cannot be a map from every bounded residual family to a bounded sampled residual family in this example. The constant residual is not claimed to be the differential of a bounded candidate on the identity tower. It is a witness about the proposed operator on the entire residual space.

Neither example contradicts cofinal invariance of coherent points or eventual transports. Those theorems concern exact compatibility. A norm on all raw defects, the retained fine detail, and the choice of a repair section are additional structures.

\begin{insight}
A compatible change of representation transports equations correctly. It transports the selected repair correctly only when the chosen residual sections also commute. The mismatch is an exact-valued operator, so it can remain invisible to a residual-only check.
\end{insight}


\chap{Multiscale residuals and the limits of reconstruction}\label{sec:cascade}
\subsection{A split coarsener retains its missing detail}
Let \(C:R\to\overline R\) and \(J:\overline R\to R\) be linear maps with \(CJ=I\). Put \(H=I-JC\). Then
\[
 r=JCr+Hr,\qquad CH=0,\qquad H^2=H.
\]
The two coordinates \((Cr,Hr)\) retain all information. The coarse coordinate alone generally does not.

The actual dyadic tower coarsener has the typed formula
\begin{equation}\label{eq:dyadic}
 (Cr)_n=\rho^Y_{2n}r_{2n+1}+r_{2n}\rho^X_{2n+1}.
\end{equation}
It is obtained by telescoping two adjacent compatibility squares. With target sections, a lift is
\[
 (Jc)_{2n}=0,\qquad (Jc)_{2n+1}=s_{2n}c_n.
\]
The section identity gives \(CJ=I\), so the missing detail is explicitly \(Hr\), not an unspecified reconstruction error. The original normed construction imposes the bounds needed to make these maps bounded on its operator family spaces.
\source{KS25}

\subsection{A whole cascade}
At scale \(j\), let \(C_j:R_j\to R_{j+1}\) and \(J_j:R_{j+1}\to R_j\) be bounded linear maps with \(C_jJ_j=I\). For \(x\in R_0\), define
\[
 x_0=x,\quad x_{j+1}=C_jx_j,\quad
 L_0=I,\quad L_N=J_0\cdots J_{N-1},
\]
\[
 h_j=x_j-J_jx_{j+1},\qquad Q_Nx=L_Nx_N,
 \qquad S_N=I-Q_N.
\]
The map \(Q_N\) restores the surviving coarse remainder to the original space. It is this restored remainder, not just \(x_N\), that controls reconstruction.

\begin{theorem}[Finite reconstruction and nested remainders]\label{KS26}
For every \(N\),
\[
 x=Q_Nx+\sum_{j<N}L_jh_j,
 \qquad S_Nx=\sum_{j<N}L_jh_j.
\]
Each \(Q_N\) is a projection, and for \(N\le M\),
\[
 Q_NQ_M=Q_MQ_N=Q_M.
\]
The vectors with all details zero are exactly the persistent subspace
\[
 \mathcal P_\infty=\{x:Q_Nx=x\text{ for every }N\}.
\]
\end{theorem}
\begin{proof}
Let $D_0=I$ and $D_N=C_{N-1}\cdots C_0:R_0\to R_N$. Thus $x_N=D_Nx$, $Q_N=L_ND_N$, and repeated use of $C_jJ_j=I$ gives
\[
 D_NL_N=I_{R_N}.
\]
This cancellation takes place between adjacent maps with matching domains. No identification of different residual spaces is needed.

Since $L_{j+1}=L_jJ_j$ and $x_{j+1}=C_jx_j$,
\[
 L_jh_j=L_j(x_j-J_jx_{j+1})=Q_jx-Q_{j+1}x.
\]
Summing from $j=0$ to $N-1$ leaves $x-Q_Nx$ because $Q_0=I$. This proves both finite reconstruction identities, including $N=0$, when the detail sum is empty.

Idempotence follows from $D_NL_N=I$:
\[
 Q_N^2=L_N(D_NL_N)D_N=Q_N.
\]
For $M\ge N$, put $C_{N,M}=C_{M-1}\cdots C_N$ and $J_{N,M}=J_N\cdots J_{M-1}$, using identities for $M=N$. Then
\[
 D_M=C_{N,M}D_N,\quad L_M=L_NJ_{N,M},\quad C_{N,M}J_{N,M}=I.
\]
The identities $D_NL_M=J_{N,M}$ and $D_ML_N=C_{N,M}$ now give
\[
 Q_NQ_M=L_NJ_{N,M}D_M=Q_M,\qquad
 Q_MQ_N=L_MC_{N,M}D_N=Q_M.
\]
Both orders have been checked, rather than inferred from one range inclusion.

If all $h_j$ vanish, the finite telescope gives $Q_Nx=x$ for every $N$. Conversely, if every $Q_Nx=x$, then $L_jh_j=Q_jx-Q_{j+1}x=0$. Apply the left inverse $D_j$ of $L_j$ to obtain $h_j=0$. The injectivity of the restoration composite is the step that recovers the original detail from its transported value.
\end{proof}
\source{KS26}

\fixedplate{05_cascade}{A split cascade retains every discarded component as a detail. The finite sum is always exact after the restored coarse remainder is included. An infinite detail expansion requires that remainder to vanish in the topology being used.}{fig:cascade}

\subsection{Convergence has its own theorem}
Call \(x\) norm recoverable when \(Q_Nx\to0\). This is equivalent to \(S_Nx\to x\), but it does not follow merely from vanishing persistent subspace. Let
\[
 \mathcal F=\bigcup_N\ker Q_N
\]
be the finite core.

\begin{theorem}[Density criterion for norm recovery]\label{KS27}
Suppose \(\sup_N\norm{Q_N}\le K<\infty\). Then the norm recoverable vectors are exactly \(\overline{\mathcal F}\). In particular all vectors are norm recoverable if and only if the finite core is dense.
\end{theorem}
\begin{proof}
The nested identities first imply that the kernels increase: if $Q_My=0$ and $N\ge M$, then $Q_Ny=Q_NQ_My=0$. Thus the finite core really is an increasing union of linear subspaces.

Suppose $Q_Nx\to0$. Put $y_N=S_Nx=x-Q_Nx$. Idempotence gives $Q_Ny_N=Q_Nx-Q_N^2x=0$, so $y_N\in\mathcal F$. Since $x-y_N=Q_Nx\to0$, these are finite-core approximants converging to $x$. Hence $x\in\overline{\mathcal F}$.

Conversely, fix $x\in\overline{\mathcal F}$ and $\epsilon>0$. If $K=0$, the common operator bound gives $\norm{Q_Nx}=0$ for every $N$, so convergence is immediate. Suppose $K>0$. Choose $y\in\mathcal F$ satisfying $\norm{x-y}<\epsilon/(K+1)$. Some $M$ has $Q_My=0$. For every $N\ge M$, nesting and the common operator bound give
\[
 \norm{Q_Nx}=\norm{Q_N(x-y)}\le K\norm{x-y}
 <\epsilon\frac{K}{K+1}\le\epsilon.
\]
The approximation $y$ is chosen before the depth $N$, giving precisely the quantifier order in norm convergence. This proves equality of the two sets. Applying it to every $x\in R_0$ gives the density criterion.
\end{proof}
\source{KS27}
There are other sufficient conditions. If the transported details are absolutely summable in the Banach space \(R_0\), their sum exists and the remaining boundary lies in \(\mathcal P_\infty\). If that persistent subspace is zero, the sum equals \(x\). Alternatively,
\[
 \norm{x-S_Nx}\le\left(\prod_{j<N}\norm{J_j}\right)\norm{x_N}
\]
gives a direct sufficient decay test. A decreasing coarse norm by itself is not enough when restoration gains grow.

\subsection{The actual dyadic result}
For the source's dyadic residual hierarchy, the persistent residual is zero. Zero detail at a scale reveals every even coordinate, while an odd coordinate is passed to a smaller index at the next scale. Strong induction on the coordinate index therefore proves that a residual with all details zero is zero.

\begin{theorem}[Dyadic determination and finite coordinate recovery]\label{KS28}
In the actual dyadic system, all detail coordinates jointly determine the residual. Moreover,
\[
 k+1<2^N\quad\Longrightarrow\quad
 (Q_Nr)_k=0,\qquad(S_Nr)_k=r_k.
\]
Thus every fixed residual coordinate is recovered exactly at a sufficiently large finite depth. The statement does not assert convergence in the supremum norm over all coordinates.
\end{theorem}
\begin{proof}
Write \(R_j\) for the residual space at scale \(j\). Its lift has the actual coordinate formula
\[
 (J_jz)_{2m}=0,\qquad (J_jz)_{2m+1}=s^{(j)}_{2m}z_m,
\]
where \(s^{(j)}_{2m}\) is the corresponding target section at that scale. The expression on the right is composition of the section with the residual operator. Its source and target are those of the odd coordinate at scale \(j\).

Set \(L_{j,0}=I_{R_j}\) and \(L_{j,N+1}=J_jL_{j+1,N}\). We prove, simultaneously for all starting scales \(j\), that
\[
 k+1<2^N\quad\Longrightarrow\quad (L_{j,N}z)_k=0
 \qquad(z\in R_{j+N}).
\]
For \(N=0\), the inequality \(k+1<1\) has no solution. Suppose the assertion holds at depth \(N\). At depth \(N+1\), an even coordinate vanishes by the displayed lift formula. For an odd coordinate \(k=2m+1\), the hypothesis gives
\[
 2(m+1)<2^{N+1},\qquad m+1<2^N.
\]
The induction hypothesis at starting scale \(j+1\) makes \((L_{j+1,N}z)_m\) zero. Composing that zero operator with \(s^{(j)}_{2m}\) gives zero at coordinate \(k\). This completes the induction without identifying residual spaces at different scales.

Now \(Q_Nr=L_{0,N}x_N\), so the support assertion applies to the restored remainder. Since \(S_Nr=r-Q_Nr\), it gives the reconstruction equation. Finally, if two residuals have identical details at every scale, the finite telescope makes their \(S_N\) reconstructions equal for every \(N\). At any fixed coordinate choose \(N\) with \(k+1<2^N\). The reconstruction equation then makes their original coordinates equal. Hence the residuals are equal.
\end{proof}
\source{KS28}
The strict inequality matters for recovery of a whole initial block. At its final coordinate the remainder can still be nonzero. There is also a sharper, nonconsecutive description of the support.

\begin{proposition}[The dyadic support pattern]\label{KS55}
For the same lifts,
\[
 2^N\nmid(k+1)\quad\Longrightarrow\quad(Q_Nr)_k=0.
\]
Thus coordinate \(k\) is certainly recovered once \(N>\nu_2(k+1)\), where \(\nu_2(m)\) is the largest power-of-two exponent dividing the positive integer \(m\). The only possible unrecovered coordinates at depth \(N\) are
\[
 2^N-1,\quad 2\cdot2^N-1,\quad3\cdot2^N-1,\quad\ldots.
\]
\end{proposition}
\begin{proof}
Prove the assertion for \(L_{j,N}\) by induction, uniformly in \(j\). At depth zero every positive integer is divisible by one, so there is nothing to prove. At depth \(N+1\), even coordinates vanish. If \(k=2m+1\), then
\[
 2^{N+1}\mid(k+1)\quad\Longleftrightarrow\quad2^N\mid(m+1).
\]
When divisibility fails, the inner coordinate in \(J_jL_{j+1,N}\) vanishes by induction. The formula for the odd lift finishes the proof. Apply the result to \(Q_Nr=L_{0,N}x_N\).
\end{proof}
\source{KS55}
This refinement is formalized by a scale-uniform induction on the actual varying-fiber residual system in \code{DyadicSupport.lean}. It asserts where a remainder \emph{may} survive, not that every allowed coordinate is nonzero. In the scalar identity-bonding model with \(C(a,b)=a+b\), the residual constantly equal to one has value \(2^N\) at every surviving position. That example proves sharpness of the uniform prefix cutoff.

\plate{16_dyadic_support}{The possible support of the restored dyadic remainder is sparse and periodic. At depth \(N\), only positions whose index plus one is divisible by \(2^N\) can remain. Open cells are guaranteed recovered. Marked cells are allowed to survive, not asserted nonzero for every residual.}{fig:dyadicsupport}

\subsection{Two boundary examples}
On \(\ell^\infty(\N,\R)\), let \(C\) remove the first coordinate and let \(J\) prepend zero. Then \(Q_Nx\) is the tail of \(x\), placed at its original coordinates with zeros before \(N\). Every fixed coordinate is eventually zero in the remainder, and \(\norm{Q_N}\le1\). Nevertheless the constant sequence \(\mathbf1\) satisfies \(\norm{Q_N\mathbf1}=1\) for every \(N\).

\begin{proposition}[The recovery domain is \(c_0\)]\label{KS29}
For this shift cascade,
\[
 S_Nx\to x\text{ in }\ell^\infty
 \quad\Longleftrightarrow\quad x_n\to0.
\]
For \(x_n=q^n\), \(0\le q<1\), the error is exactly \(q^N\).
\end{proposition}
\begin{proof}
The error norm is \(\sup_{n\ge N}|x_n|\). It tends to zero precisely when the sequence tends to zero. For a geometric sequence its largest tail value is the first one, \(q^N\).
\end{proof}
\source{KS29}
As a second example, take scalar spaces with \(C=\tfrac12I\), \(J=2I\). The coarse value \(x_N=2^{-N}x\) tends to zero, but restoration gives \(Q_Nx=x\). Every detail is zero and the entire input is persistent. The example belongs to this scalar cascade, not to the actual dyadic hierarchy, whose persistent subspace was proved to vanish.

\begin{insight}
Retained packets make omission explicit. They do not guarantee convergence. Coordinate recovery, bounded reconstruction, summability, and attainment remain different mathematical properties.
\end{insight}
